\chapter{Rubrica Sugerida de Avaliação}

Esta rubrica é uma sugestão de apoio. Cada curso pode definir critérios, pesos e instrumentos próprios. Use-a também de forma formativa para identificar fragilidades e discutir prioridades, sem transformar o julgamento acadêmico em checklist mecânico.

\begin{longtable}{p{0.22\linewidth}p{0.30\linewidth}p{0.30\linewidth}}
\toprule
Critério & Desempenho satisfatório & Pontos de atenção \\
\midrule
\endfirsthead
\toprule
Critério & Desempenho satisfatório & Pontos de atenção \\
\midrule
\endhead
Tema e problema & Tema delimitado, problema claro e relevante. & Tema amplo, pergunta vaga ou recorte inviável. \\
Objetivos & Objetivos coerentes e verificáveis. & Objetivos genéricos ou desconectados do método. \\
Justificativa & Relevância demonstrada com argumentos e fontes. & Justificativa opinativa ou superficial. \\
Referencial teórico & Conceitos articulados ao problema. & Revisão descritiva, desatualizada ou desconexa. \\
Metodologia & Procedimentos claros e adequados. & Método genérico, incompleto ou incompatível. \\
Análise & Resultados interpretados com rigor. & Dados apenas descritos ou sem diálogo teórico. \\
Conclusão & Responde ao problema e reconhece limites. & Repete introdução ou traz informações novas. \\
Escrita & Texto claro, coeso e acadêmico. & Ambiguidade, informalidade ou problemas recorrentes de redação. \\
Citações e referências & Fontes adequadas e padronizadas. & Ausências, inconsistências ou uso inadequado de citações. \\
Apresentação formal & Trabalho conforme padrão institucional. & Paginação, sumário, listas ou elementos obrigatórios com falhas. \\
Defesa oral & Domínio do percurso e respostas consistentes. & Desconhecimento do próprio texto ou respostas improvisadas. \\
\bottomrule
\end{longtable}
